\label{sec:GeneralLimitations}
\subsection{AES and Label-Size}\label{sub:AESSizeLimitations}
The given code by Vadapalli et al. is optimized for 64-bit values with 128-bit labels.
These labels are represented by Intel's \code{\_\_m128i} SIMD-datatype (\ref{subsub:implementationAesni}).
Keys for AES are represented using an array holding eleven \code{\_\_m128i} values, which means they are of type \code{\_\_m128i[11]}.
Eleven values are necessary since AES for 128-bit values takes place in eleven rounds, thus eleven keys.
Adjusting AES for 256 or 512 bits requires more implementations using different inputs and keys, since a higher number of rounds is needed.

AES as used in the implementation keeps its secret-premises when concatenated.
This means, executing AES multiple times and concatenating the results into a single word is still considered secure.
Also, implementing arbitrarily scalable AES far exceeds this research's efforts.
Due to this reasons, the author chose to scale AES by concatenating 128-bit inputs.

This results in the requirements that all inputs, or the labels of given inputs, must be a multiple of 128 bits.
These inputs then can be splitted into 128-bit blocks, encrypted or decrypted using the 128-bit AES implementation, and finally concatenated to produce a single result.
This result is then the same length as the input, which is the equivalent of executing AES on the complete input instead of the divided block structure.

\subsection{WIDTH}\label{sub:WIDTH}
